\section{Mathematical Foundation}
\label{sec:math}

\subsection{Domains}

Let $\mathcal{T}$ denote the set of all tensor types (rank, shape, data
type); $\mathcal{O}$ the set of tensor operations, i.e.\ partial
functions from input types to output types; $\mathcal{B}$ the set of
hardware backends, each characterized by its supported operations,
data types, memory capacity, and alignment requirements; and
$\mathcal{L}$ the set of memory layouts, i.e.\ mappings from logical
tensor indices to physical offsets. A \emph{tensor computation graph}
is a DAG $\mathcal{G} = (V, E)$ whose vertices are tensor operations
and whose edges are data dependencies.

\subsection{Execution Plans}

\begin{definition}[Execution plan]
An execution plan for $\mathcal{G}$ is a tuple
$P = (\sigma, \beta, \mathcal{L}_{\mathrm{map}}, \mathbf{m})$, where
$\sigma = (o_1, \ldots, o_k)$ is a topological ordering of the
operations of $\mathcal{G}$; $\beta \in \mathcal{B}$ is the target
backend; $\mathcal{L}_{\mathrm{map}}: V \to \mathcal{L}$ assigns a
memory layout to each intermediate tensor; and
$\mathbf{m} \in \R^d$ is the plan's metric vector.
\end{definition}

The set of all plans is $\mathcal{P}$; a finite candidate subset is
$\mathcal{P}_{\text{cand}} \subseteq \mathcal{P}$.

\subsection{Constraints and Validation}

\begin{definition}[Constraint]
A constraint is a \emph{total} function
$C: \mathcal{P} \to \{0, 1\}$. Plan $P$ satisfies $C$ iff $C(P) = 1$.
\end{definition}

Totality guarantees that every plan receives a definitive pass/fail
determination---there are no undefined cases, which is essential for
the correctness results of Section~\ref{sec:correctness}.

\begin{definition}[Validator]
A validator over constraints $C_1, \ldots, C_k$ is the product
\begin{equation}
V(P) \;=\; \prod_{i=1}^{k} C_i(P).
\label{eq:validator}
\end{equation}
\end{definition}

$V(P) = 1$ iff every constraint accepts $P$; a single violation
invalidates the plan. This strict conjunctive semantics encodes the
principle that correctness is all-or-nothing.

\begin{definition}[Valid execution set]
$\mathcal{P}_{\text{valid}} =
 \{P \in \mathcal{P}_{\text{cand}} \mid V(P) = 1\}$.
\end{definition}

$\mathcal{P}_{\text{valid}}$ is always well defined (possibly empty),
and it is the \emph{only} domain over which the \gate{} optimizer
operates.

\subsection{Metrics, Weights, and Cost}

\gate{} uses a five-dimensional metric vector
$\mathbf{m}(P) = (m_1, \ldots, m_5)(P)$:

\begin{center}
\small
\begin{tabular}{@{}llll@{}}
\toprule
Dimension & Symbol & Description & Unit \\
\midrule
Latency & $m_1$ & estimated total execution time & ms \\
Peak memory & $m_2$ & maximum live allocation & MB \\
Sync overhead & $m_3$ & kernel launches, transfers, barriers & ms \\
Energy & $m_4$ & total energy consumed & mJ \\
Numerical error & $m_5$ & cumulative $L_2$ error estimate & --- \\
\bottomrule
\end{tabular}
\end{center}

All components are non-negative and are estimated analytically from the
plan's structure, so cost evaluation itself adds no execution overhead.

\begin{definition}[Weight vector]
$\mathbf{w} \in \R^d_{\geq 0}$ with $\sum_{i=1}^d w_i = 1$; the set of
valid weight vectors is the standard simplex $\Delta^{d-1}$.
\end{definition}

\begin{definition}[Cost function]
$\mathcal{C}(P, \mathbf{w}) = \mathbf{w}^{\top} \mathbf{m}(P)
 = \sum_{i=1}^{d} w_i \, m_i(P)$.
\end{definition}

\begin{lemma}[Monotonicity]
\label{lem:monotonicity}
For fixed $\mathbf{w}$ with $w_i > 0$, $\mathcal{C}$ is strictly
increasing in each metric component: if $P$ and $P'$ differ only in
dimension $i$ with $m_i(P) > m_i(P')$, then
$\mathcal{C}(P, \mathbf{w}) > \mathcal{C}(P', \mathbf{w})$.
\end{lemma}

\begin{proof}
Immediate from linearity:
$\mathcal{C}(P,\mathbf{w}) - \mathcal{C}(P',\mathbf{w})
 = w_i \,(m_i(P) - m_i(P')) > 0$.
\end{proof}

\subsection{Execution Policies}

\begin{definition}[Execution policy]
A policy is a mapping $\pi: \mathcal{K} \to \Delta^{d-1}$ from a
deployment context $\kappa$ to a weight vector.
\end{definition}

\gate{} ships five presets: Latency $(0.6,0.1,0.1,0.1,0.1)$, Memory
$(0.1,0.6,0.1,0.1,0.1)$, Energy $(0.1,0.1,0.1,0.6,0.1)$, Balanced
$(0.2,0.2,0.2,0.2,0.2)$, and Custom (user-supplied).

\begin{theorem}[Policy independence of validity]
\label{thm:policy-independence}
For any $\mathbf{w}, \mathbf{w}' \in \Delta^{d-1}$,
$\mathcal{P}_{\text{valid}}(\mathbf{w}) =
 \mathcal{P}_{\text{valid}}(\mathbf{w}')$.
\end{theorem}

\begin{proof}
$\mathcal{P}_{\text{valid}}$ is defined by $V$ alone
(Eq.~\ref{eq:validator}); no constraint reads $\mathbf{w}$. Hence the
valid set is invariant under any change of weights.
\end{proof}

The policy affects only \emph{which valid plan is selected}, never
\emph{whether the selected plan is valid}.

\subsection{The Selection Criterion}

The central operation of \gate{} is
\begin{equation}
P^{*} \;=\; \argmin_{P \in \mathcal{P}_{\text{valid}}}
            \mathcal{C}(P, \mathbf{w})
      \;=\; \argmin_{P \in \mathcal{P}_{\text{cand}}}
            \bigl\{\, \mathcal{C}(P, \mathbf{w}) \;\big|\; V(P) = 1 \,\bigr\}.
\label{eq:selection}
\end{equation}

The side condition $V(P) = 1$ is the mathematical expression of
``constraint before optimization.''

\subsection{Empty Valid Set Semantics}

When $\mathcal{P}_{\text{valid}} = \emptyset$ the argmin is undefined,
and \gate{} raises an explicit \texttt{ConstraintViolation} error. The
dispatch function is the partial function
\begin{equation}
\Phi(\mathcal{P}_{\text{cand}}, V, \mathbf{w}) =
\begin{cases}
\displaystyle P^{*} = \argmin_{P \in \mathcal{P}_{\text{valid}}}
  \mathcal{C}(P, \mathbf{w})
  & \text{if } \mathcal{P}_{\text{valid}} \neq \emptyset, \\[1.2em]
\bot \;\; (\texttt{ConstraintViolation})
  & \text{if } \mathcal{P}_{\text{valid}} = \emptyset.
\end{cases}
\end{equation}

\gate{} either returns a valid, cost-minimal plan or raises an explicit
error; it never returns an invalid plan.
